% 作者题证的编辑转录副本，不是独立下载的上游原件。
% 来源：https://terrytao.wordpress.com/2011/09/08/254a-notes-2-building-lie-structure-from-representations-and-metrics/

\begin{theorem}[Cartan's closed-subgroup theorem; Tao's Theorem 2]
A closed subgroup $H$ of a finite-dimensional real Lie group $G$ is a
smooth embedded submanifold of $G$, and hence is a Lie group with its
subspace topology.
\end{theorem}
\paragraph{One-parameter subgroups.}
A one-parameter subgroup is a continuous homomorphism $\R\to G$.
For a Lie group these are exactly the maps $t\mapsto\exp(tX)$,
$X\in\mathfrak g$. Thus the one-parameter subgroups of $H$ identify with
\[
 \mathfrak h=\{X\in\mathfrak g:\exp(tX)\in H\text{ for all }t\in\R\}.
\]
This set contains zero and is closed under real scalar multiplication.
If $X,Y\in\mathfrak h$, the Lie product formula gives
\[
 \exp(t(X+Y))=
 \lim_{n\to\infty}\bigl(\exp(tX/2^n)\exp(tY/2^n)\bigr)^{2^n}.
\]
Every term lies in $H$, so its limit lies in $H$ by closedness.
Consequently $\mathfrak h$ is a linear subspace of $\mathfrak g$.

\begin{lemma}[Nontrivial tangent direction; Tao's Lemma 10]
If the identity is not isolated in $H$, then $\mathfrak h\ne\{0\}$.
\end{lemma}
\begin{proof}
Choose $h_n\in H\setminus\{e\}$ with $h_n\to e$.
The exponential map is a local diffeomorphism near zero, so for large $n$
write $h_n=\exp(X_n)$ with $X_n\ne0$, $X_n\to0$.
Fix a norm on the finite-dimensional vector space $\mathfrak g$.
By compactness of its unit sphere, after passing to a subsequence
\[
 X_n/\|X_n\|\longrightarrow\omega,\qquad\|\omega\|=1.
\]
For $t>0$,
$\lfloor t/\|X_n\|\rfloor X_n\to t\omega$ and therefore
\[
 h_n^{\lfloor t/\|X_n\|\rfloor}
 =\exp\bigl(\lfloor t/\|X_n\|\rfloor X_n\bigr)
 \longrightarrow\exp(t\omega).
\]
The powers belong to $H$; their limits do too. Taking inverses also covers
$t<0$, so $\omega\in\mathfrak h$ and $\mathfrak h$ is nontrivial.
\end{proof}

\begin{lemma}[Local exponential chart; Tao's Lemma 11]
There are identity and zero neighbourhoods $U\subset H$ and
$V\subset\mathfrak h$ for which $\exp:V\to U$ is a homeomorphism.
\end{lemma}
\begin{proof}
Take a small neighbourhood $V$ of zero in $\mathfrak h$ on which the
exponential map is injective and is a homeomorphism onto its image.
If $\exp(V)$ contains a neighbourhood of the identity in $H$, the claim
follows after restricting neighbourhoods. Otherwise there are
$h_n\in H\setminus\exp(V)$ with $h_n\to e$. Write
$h_n=\exp(X_n)$, where $X_n\to0$ in $\mathfrak g$.

Choose a vector-space complement $\mathfrak k$ of $\mathfrak h$:
\[
 \mathfrak g=\mathfrak h\oplus\mathfrak k.
\]
No Lie-algebra condition is imposed on $\mathfrak k$. By the inverse
function theorem, the smooth map
\[
 \mathfrak h\times\mathfrak k\longrightarrow G,
 \qquad(Y,Z)\longmapsto\exp(Y)\exp(Z)
\]
is a local diffeomorphism at $(0,0)$. For all large $n$ we may write
\[
 h_n=\exp(Y_n)\exp(Z_n),\qquad
 Y_n\in\mathfrak h,\quad Z_n\in\mathfrak k,\quad Y_n,Z_n\to0.
\]
Since $h_n\notin\exp(V)$ and $Y_n$ is eventually in $V$, we have $Z_n\ne0$.
Choose a norm on $\mathfrak k$ and pass to a subsequence such that
$Z_n/\|Z_n\|\to\omega$ with $\|\omega\|=1$. In particular,
$\omega\in\mathfrak k\setminus\mathfrak h$.

But $\exp(Z_n)=\exp(-Y_n)h_n\in H$. The integer-power argument in the
previous lemma now implies $\exp(t\omega)\in H$ for every $t\in\R$.
Thus $\omega\in\mathfrak h$, a contradiction.
\end{proof}

\begin{proof}[Conclusion of Cartan's theorem]
The preceding lemma shows that $H$ agrees locally near the identity with
$\exp(V)$, where $V$ is open in the linear subspace $\mathfrak h$.
The ambient exponential chart is smooth with smooth inverse, so this is
a smooth embedded submanifold near the identity. Translation by elements
of $H$ gives such charts at every point of $H$. The group operations are
restrictions of the smooth operations of $G$, hence are smooth in this
submanifold structure. Therefore $H$ is a Lie group.
\end{proof}
